\scanpage{25}
\mspage{23}
% archive sheet number: 24
% hand: Grothendieck
% sheet_kind: text
% element: text | top, lines 1–2 | G^(i+1) = [G^(i),G]
$G^{(i+1)}=[G^{(i)},G]$) tel que $\operatorname{Gr}(G^\circ)$ soit un groupe vectoriel. On en conclut, comme
% element: display | upper, line 3 | H^i(Gamma,V)=0
\[
H^{\unsure{i|\nu}}(\Gamma,V)=0\qquad\text{pour }i\geq0
\]
% element: text | upper, lines 4–7 | et tout vectoriel V
et $\unsure{\text{tt}|\forall}$ vectoriel $V$ sur lequel $\Gamma$ opère (car $V$ est uniquement divisible\ldots), que \uline{l'extension} $G$ de $\Gamma$ par $G^\circ$ \uline{splitte de façon unique en un} produit semi-direct
% element: display | upper-middle, line 8 | G isom Gamma.G^0
\[
G\simeq\Gamma.G^\circ.
\tag*{(5.5.1)}
\]
% element: text | middle, lines 9–11 | Ceci montre que
Ceci \deleted{\unsure{permet de voir que}}\inserted{montre que $\Gamma$ opère de façon naturelle} sur $E$, \deleted{[et permet de retrouver le fait que $\forall g\in G(R)$, \unsure{i.e.} $g$ unipotent $\Rightarrow g\in G^\circ(R)$]}
% element: margin note | left, below cancelled paragraph | (5.6), 5.6
\margin{\deleted{(5.6)}\quad\deleted{5.6}}
% element: text | middle, unfinished last line | Pour t.
\unsure{Pour t.}
